\origpage{279}{GRAPH MINORS. VIII}
\noindent \textit{$k$-train} over $\mathcal{A}$ if it has a path-decomposition $(G_i, Y_{i-1}, Y_i)$ ($1 \le i \le r$) over $\mathcal{A}$ such that

\hspace*{12pt}(i) for $1 \le i < r$, $|Y_i| = k$

\hspace*{12pt}(ii) if $r \ge 3$ there are $k$ mutually vertex-disjoint paths of $G$ between $Y_1$ and $Y_{r-1}$.

If $\mathcal{A}$ is a class of grafts, we denote by $\mathcal{A}^*$ the class of all grafts isomorphic (in the natural sense) to a member of $\mathcal{A}$. Evidently if $\mathcal{A}$ is well-rooted then so is $\mathcal{A}^*$.

\stmt
(14.3) \textit{Every $k$-train over $\mathcal{A}$ belongs to $W(\mathcal{A}^* \oplus \mathcal{B}^* \oplus \mathcal{A}^*) \cup W(\mathcal{A}^* \oplus \mathcal{A}^*) \cup \mathcal{A}$ where $\mathcal{B}$ is the class of all $k$-chains over $\mathcal{A}$.}

\textit{Proof.} Let $(G, T)$ be a $k$-train over $\mathcal{A}$, and let $(G_i, Y_{i-1}, Y_i)$ ($1 \le i \le r$) be the corresponding path-decomposition. If $r = 1$ then $(G, T) \in \mathcal{A}$. If $r = 2$ then $(G, T) \in W(\mathcal{A}^* \oplus \mathcal{A}^*)$. If $r \ge 3$ then
\begin{equation*}
(G, T) \in W(\mathcal{A}^* \oplus \mathcal{B}^* \oplus \mathcal{A}^*).
\end{equation*}
In each case the theorem is true, as required. $\square$

If $\mathcal{A}$ is a class of grafts, $L^k(\mathcal{A})$ denotes the class of $k$-trains over $\mathcal{A}$, and $C^k(\mathcal{A})$ denotes $L^0(L^1(L^2(...(L^k(\mathcal{A}))...)))$.

\stmt
(14.4) \textit{If $\mathcal{A}$ is well-rooted then so is $L^k(\mathcal{A})$.}

This follows immediately from (14.2), (13.1), (13.3), (14.3), and (4.2).

We say that a graft $(G, T)$ is a \textit{concatenation over $\mathcal{A}$} if $(G, T) \in C^k(\mathcal{A})$ for some $k \ge 0$.

\stmt
(14.5) \textit{If $\mathcal{A}$ is well-rooted, then the class of all concatenations over $\mathcal{A}$ is well-rooted.}

\textit{Proof.} By (4.1), $\mathcal{A}$ has some index $k_0$. Now for $k > k_0$, $L^k(\mathcal{A}) = \mathcal{A}$, and so every concatenation over $\mathcal{A}$ is an element of $C^{k_0}(\mathcal{A})$. But this is well-rooted, from (14.4). The result follows. $\square$

\sectionhead{15. \textsc{Tubes}}

A \textit{cylinder} is a surface homeomorphic to $\Sigma(0, 0, 2)$. Our next objective is to prove (6.1) when $\Sigma$ is a cylinder.

There are two kinds of $O$-arc in a cylinder; those that are null-homotopic and bound a disc, and those that wind once around $\Sigma$. A \textit{hoop} is an $O$-arc of the second kind. If $F$ is a hoop then $\Sigma - F$ is naturally divided into two
